\documentclass[11pt, letterpaper, includehead]{article}

%%%%%%%%%%%%%%%%%%%%% Pre-document %%%%%%%%%%%%%%%%%%%%%
\usepackage{fancyhdr}  % Allow for headers
\usepackage{graphicx}  % Allow for figures 
\usepackage{float}     % Allow for figure inserted in specified location
\usepackage{xcolor}

\setlength{\parindent}{0pt} % Remove auto paragraph indents

% Get rid of those big ass margins
\usepackage[margin=1in]{geometry}

\begin{document}
  %%%%%%%%%%%%%%%%%%%%% Title Page %%%%%%%%%%%%%%%%%%%%%
  \begin{titlepage} 
    \begin{center}
      \Huge{\textbf{Lab 1}}\\
      \Huge{Motion diagrams}
      \vfill
      \begin{figure}[H] % H makes the figure insert at the position in the document
        \centering 
        \includegraphics[width=6cm]{r_cubed_logo.png}
      \end{figure}
      % \large{\textbf{Rectangle Repulsed Researchers}}\\
      \large{Julian Barossi, Liam Gilligan, Stephanie L'Heureux}\\
      \vspace{0.5cm}
      \normalsize
      \today
    \end{center}
  \end{titlepage}

  %%%%%%%%%%%%%%%%%%%%% TABLE OF CONTENTS %%%%%%%%%%%%%%%%%%%%%
  \tableofcontents
  \pagebreak % Move to next page


  % Add a nice fancy header
  \pagestyle{fancy}
  \fancyhead{}
  \fancyhead[C]{\textbf{Lab 1:} Motion diagrams}

  
  \setcounter{section}{2} % Make numbers start at 3
  \section{Matching Position vs. Time and Velocity vs. Time Graphs}

  %%%%%%%%%%%%%%%%%%%%% GRAPH 2 %%%%%%%%%%%%%%%%%%%%%
  \subsection{Graph 1}

  \begin{figure}[H] % H makes the figure insert at the position in the document
    \centering 
    \includegraphics[width=\linewidth]{graph_1.png}
  \end{figure}

  \subsubsection{Graph analysis}
  \emph{Graph 1} describes the position vs. time graph of a collision.
  The graph directly conveys that initially the object's position from 
  the sensor increases, then at $t\approx2.4s$ it's position begins decreasing. 
  The derivative of a position vs. time graph gives the velocity. From the slope 
  of the graph, one can conclude that the velocity over the interval $t\epsilon[0s, 2.4s)$
  is constant and positive. Then from $t \epsilon (2.4s, 6s]$ the velocity is constant 
  and negative. With this information, one can deduce the object of interest was given 
  an initial velocity and began moving in the positive direction, then it encountered 
  a barricade and reversed direction.
  $$\forall \, t \in [0s, 2.4s); \, v(t) = \frac{d}{dt}[s(t)] \approx \frac{0.95m - 0.6m}{2s - 1s} = 0.35 \, m/s$$
  $$\forall \, t \in (2.4s, 6s]; \, v(t) = \frac{d}{dt}[s(t)] \approx \frac{0.95m - 0.7m}{3s - 4s} = -0.25 \, m/s$$

  \subsubsection{Setup}
  To reproduce \emph{Graph 1}, we utilized a level track with a magnetic bumper mounted $60cm$ 
  from the \emph{PASCO Motion Sensor II}. With the cart positioned near the sensor, and we 
  pushed and released it towards the bumper. Upon impact with the bumper, it reversed 
  direction. The motion captured by the sensor was displayed using the \emph{PASCO Universal 
  850 Interface}. We consider friction and drag negligible. \\

  %%%%%%%%%%%%%%%%%%%%% GRAPH 2 %%%%%%%%%%%%%%%%%%%%%
  \subsection{Graph 2}

  % Centered figure in text location
  \begin{figure}[H] % H makes the figure insert at the position in the document
    \centering 
    \includegraphics[width=\linewidth]{graph_2.png}
  \end{figure}

  \subsubsection{Graph analysis}
  \emph{Graph 2} depicts an object whose speed is increasing at a constant rate. This deduction is 
  justified by graphical analysis. The graph shows that position increases over the interval 
  $t\epsilon[0s, 4.3s)$ at a constant rate, resembling a parabolic curve. The derivative of the 
  position vs. time graph yields velocity, which for this plot is increasing and positive. One 
  possible senerio that would recreate the graph is an object which accelerates down an inclined plane.\\
  
  \subsubsection{Setup}
  To replicate \emph{Graph 2}, we positioned a track such that it had a downward slope with the 
  \emph{PASCO Motion Sensor II} at the bottom. We released the cart from the top of the incline 
  and let it slide towards the sensor. The \emph{PASCO Universal 850 Interface} was used 
  to plot motion data collected with sonar. We consider friction and drag negligible in 
  this application.
  
  %%%%%%%%%%%%%%%%%%%%% GRAPH 3 %%%%%%%%%%%%%%%%%%%%%
  \subsection{Graph 3}

  \begin{figure}[H] % H makes the figure insert at the position in the document
    \centering 
    \includegraphics[width=\linewidth]{graph_3.png}
  \end{figure}

  \subsubsection{Graph analysis}
  \emph{Graph 3} illustrates the position vs. time curve of an object in freefall. 
  We know this because the graph shows the object's position starts from a height of 
  $1.2m$. On the interval $t\epsilon[2.3s, 2.9s)$, the motion of the object follows a 
  parabolic curve, which is typical for freefall. To further support this assumption,
  we can examine the derivative of the position vs. time graph, which is velocity. 
  The slope of the tangent line or the velocity over this interval is negative, 
  which means the object is moving in the negative direction, or falling. At $t = 2.9$ 
  the curve becomes a straight line, this is where the object impacts the ground and rests. 
  Through this analysis, conclude the object was released from a height of $1.2m$ and 
  fell to the ground. Upon impact, the object came to a rest.\\
  
  \subsubsection{Setup}
  To best replicate the motion depicted by the graph, we positioned the \emph{PASCO Motion Sensor II} 
  on the ground and tilted it so the sensor was facing upward. To protect the equipment, 
  a wire basket was placed overtop. For the object, we chose a large beach ball. 
  The beach ball was droped above the sensor and its motion was plotted by the 
  \emph{PASCO Universal 850 Interface}.

  %%%%%%%%%%%%%%%%%%%%% GRAPH 4 %%%%%%%%%%%%%%%%%%%%%
  \subsection{Graph 4}

  \begin{figure}[H] % H makes the figure insert at the position in the document
    \centering 
    \includegraphics[width=\linewidth]{graph_4.png}
  \end{figure}

  \subsubsection{Graph analysis}
  \emph{Graph 4} depicts an object moving down an incline. It comes in contact with somethinhg at the
  bottom and changes direction. As it changes direction it slows down. Then at $t=2.6s$ its velocity 
  becomes negaiive but the acceleration staus positive where it is slowing down\\
  $$\forall \, t \in [0s, 2.6s); \, a(t) = \frac{d}{dt}[v(t)] \approx \frac{0.6 \, m/s - 0.2 \, m/s}{2s - 1s} = 0.4 \, m/s^2$$
  $$\forall \, t \in (2.6s, 4.1s]; \, a(t) = \frac{d}{dt}[v(t)] \approx \frac{-0.5 \, m/s -(-0.4 \, m/s)}{2.6s - 2.8s} = 0.5 \, m/s^2$$ % Should be 0.4
  
  \subsubsection{Setup}
  Our setup consisted of a the sensor positioned at the biotton of an incline plane. We released the cart from the top 

  %%%%%%%%%%%%%%%%%%%%% GRAPH 5 %%%%%%%%%%%%%%%%%%%%%
  \subsection{Graph 5}

  \begin{figure}[H] % H makes the figure insert at the position in the document
    \centering 
    \includegraphics[width=\linewidth]{graph_5.png}
  \end{figure}

  \subsubsection{Graph Analysis}
  \emph{Graph 5} represents the velocity vs. time graph of an object that starts from rest and 
  then subsequently experiences a force. This assumption is supported by the fact that 
  the initial velocity recorded is a constant $0\,m/s$, indicating the object is originally 
  not in motion. At $t=2.4$, velocity increases rapidly and continues to escalate over the 
  interval $t\epsilon(2.4s, 2.8s)$.
  Acceleration is the derivative of the velocity vs. time graph. The slope over this interval 
  is constant, so acceleration is constant. Furthermore, the positive acceleration and positive 
  velocity indicate the object is speeding up.
  Acceleration can be defined in terms of Newton's Second Law as $$\sum \vec{F} = m \vec{a} \Rightarrow \vec{a} = \frac{\sum\vec{F}}{m} $$
  This means for acceleration to occur, there must be a net force. Therefore, we can conclude over the interval, 
  a net force is applied to the object. At $t=2.8s$, the velocity becomes constant once more which means 
  $\vec{a}=0\,m/s^2$ and $\vec{F_{net}=0N}$ so the object moves at a steady speed. 
  The motion depicted by the graph likely describes an object initially at rest and then 
  experienced a force over the interval $t\epsilon(2.4s, 2.8s)$ causing it to speed up. 
  After that, it maintained a constant velocity.
  $$\forall \, t \in (2.4s, 2.8s); \, a(t) = \frac{d}{dt}[v(t)] \approx \frac{0.1 \, m/s - 0.24 \, m/s}{2.6s - 2.7s} = 1.4 \, m/s^2$$

  \subsubsection{Setup}
  Our setup consisted of the \emph{PASCO Motion Sensor II} mounted to the end of a level track.
  We began data collection with the \emph{PASCO Universal 850 Interface}, then pushed the cart 
  away from the sensor and then released it. We consider friction and drag negligible.


  %%%%%%%%%%%%%%%%%%%%% GRAPH 6 %%%%%%%%%%%%%%%%%%%%%
  \subsection{Graph 6}

  \begin{figure}[H] % H makes the figure insert at the position in the document
    \centering 
    \includegraphics[width=\linewidth]{graph_6.png}
  \end{figure}

  \subsubsection{Graph Analysis}
  The graph above depicts an object being sped up rapidly before decelerating
  at a constant rate.
  $$\forall \, t \in (2, 2.5s); \, a(t) = \frac{d}{dt}[v(t)] \approx \frac{0.08 \, m/s - 0.38 \, m/s}{2.1s - 2.3s} = 1.5 \, m/s^2$$
  $$\forall \, t \in (2.6s, 4.1s]; \, a(t) = \frac{d}{dt}[v(t)] \approx \frac{0.6 \, m/s -(-0.4 \, m/s)}{2.8s - 4.8s} = -0.5 \, m/s^2$$ 
  % Graph 6: ~ −0.53m/s/s
  \subsubsection{Setup}
  To recreate the motion described by the graph, we decided to use an inclined plane.
  The graph initially shows a positive acceleration, but later a constant negative acceleration.
  To recreate this on the inclined plane, we pushed the cart up, then let
  gravity slow it down and accelerate it back down the incline.

  %%%%%%%%%%%%%%%%%%%%% GRAPH 7 %%%%%%%%%%%%%%%%%%%%%

  \subsection{Graph 7}

  \begin{figure}[H] % H makes the figure insert at the position in the document
    \centering 
    \includegraphics[width=\linewidth]{graph_7.png}
  \end{figure}

  \subsubsection{Graph Analysis}
  Graph 7 depicts some sort of oscillating motion. Because the graph is 
  sinusoidal, we know that the graph's position and acceleration will look
  just like the velocity, just shifted by some amount. This means that the 
  position is osciallting, as is the velocity and the acceleration.
  \subsubsection{Setup}
  In order to achieve the oscillating motion displayed, we decided to attach
  a weight to a spring. If we stretched the spring by some amount, then
  let go as we started recording the velocity, we knew that we could capture
  the same motion described by the graph.




  % \centerline{\textbf{Graph 1}}\\
  % % Any kind of straight lines, calculate a slope. D
  % \textbf{Graph 2}\\
  % \textbf{Graph 3}\\
  % \textbf{Graph 4}\\
  % \textbf{Graph 5}\\
  % \textbf{Graph 6}\\
  % \textbf{Graph 7}\\

  % \centerline{}\includegraphics{graph_1.png}}

\end{document}

